\documentclass[a4paper,10pt]{article} \usepackage{anysize}
\marginsize{2cm}{2cm}{1cm}{1cm}
%\textwidth 6.0in \textheight = 664pt
\usepackage{xltxtra}
\usepackage{xunicode} \usepackage{graphicx}
\usepackage{color} \usepackage{xgreek} \usepackage{fancyvrb}
\usepackage{minted}
\usepackage{listings}
\usepackage{enumitem} \usepackage{framed} \usepackage{relsize}
\usepackage{float} \setmainfont[Mapping=TeX-text]{DejaVu Sans}
\begin{document}

\begin{titlepage}
\begin{center}
\begin{figure}[t]
     \includegraphics[scale=0.7]{title/ntua_logo}
\end{figure}
\begin{LARGE}\textbf{ΕΘΝΙΚΟ ΜΕΤΣΟΒΙΟ ΠΟΛΥΤΕΧΝΕΙΟ\\}\end{LARGE}
\vspace{2cm}
\begin{Large}
ΣΧΟΛΗ ΗΜ\&ΜΥ\\
Τεχνητή Νοημοσύνη\\
1\textsuperscript{η} Άσκηση\\
Ακ. έτος 2011-2012\\
\end{Large}
\vspace{5cm}
\vspace{1cm}
\begin{tabular}{l r}
\Large{Γερακάρης Βασίλης}&
\large{Α.Μ.: 03108092}\\
\Large{Λύρας Γρηγόρης}&
\large{Α.Μ.: 03109687}\\
\end{tabular}\\
\vspace{5cm}

\vfill
\large\today\\
\end{center}
\end{titlepage}



\section*{Υλοποίηση αλγορίθμου A*} \setcounter{section}{1}
Ο πηγαίος κώδικας της main.c που κληθήκαμε να γράψουμε ήταν ο εξής:

Στη συνέχεια δημιουργήσαμε το makefile για  τη μεταγλώττιση του προγράμματος
με τα εξής περιεχόμενα:


Τρέχοντας στο shell την εντολή make έχουμε την παρακάτω έξοδο

\section*{Απαντήσεις στις θεωρητικές ερωτήσεις}
Kakamaka
\end{document}
